\Needspace{6\baselineskip}
\begin{center}\hypertarget{native-262-article189}{}\pdfbookmark[2]{189.}{article-189}189.\end{center}
Methodus haecce a praecedente in eo tantum differt, quod duobus locis
geocentricis exacte, \'{a}c dein reliquis quam exactissime satisfit, dum secundum
methodum alteram observatio nulla reliquis praefertur, sed errores quantum
fieri potest inter omnes distribuuntur\rlap{\raisebox{1.4ex}{\rule{.6em}{.35pt}}}. Methodus art. praec. itaque priori
eatenus tantum postponenda erit, quatenus locis $L,L'$ aliquam errorum partem
recipientibus errores in locis reliquis notabiliter diminuere licet: attamen plerumque
per idoneam electionem observationum $L,L'$ facile caveri potest, ne
haec differentia magni momenti evadere possit. Operam scilicet dare \SourceLeaf{263}oportebit,
ut pro $L,L'$ tales observationes adoptentur, quae non solum exquisita
praecisione gaudeant, sed ita quoque comparatae sint, ut elementa ex ipsis
distantiisque derivata a variationibus parvis ipsarum positionum geocentricarum
non nimis afficiantur. Parum prudenter itaque ageres, si observationes parvo
temporis intervallo ab invicem distantes eligeres, talesve, quibus loci heliocentrici
proxime oppositi vel coincidentes responderent.
\begin{center}\rule{3cm}{.4pt}\end{center}
